Теорема. Для любого выпуклого четырёхугольника (трапеции) $A B C D$ с диагоналями $A C$ и $B D$ , если построить параллелограмм на сторонах $A B$ и $B C$, получив параллелограмм $A B C E$, и соединить точки $D$ и $E$ отрезком $D E$, то сумма квадратов четырёх сторон четырёхугольника $A B^2+B C^2+C D^2+D A^2$ будет больше суммы квадратов диагоналей $A C^2+B D^2$ на квадрат отрезка $D E$, то есть

$$
A B^2+B C^2+C D^2+D A^2=A C^2+B D^2+D E^2 .
$$


Доказательство. Сначала достроим три точки $A, B, C$ до параллелограмма $A B C E$, как предложено в условии теоремы, и проведём диагональ $B E$. Затем построим точку $F$ так, чтобы $C F$ была параллельна $A D$, а $B F$ - параллельна $E D$. Поскольку $B C=A E$, получаем, что треугольники $C B F$ и $A E D$ равны.

Теперь проведём линии $A F, D F$ и $E F$ и рассмотрим два параллелограмма $A D C F$ и $B D E F$ с диагоналями $A C, D F$ и $B E, D F$ соответственно. Эйлер ссылается на «свойство параллелограммов», согласно которому в $A D C F$ выполняется

$$
2 A D^2+2 C D^2=A C^2+D F^2
$$


а в $B D E F$ выполняется

$$
2 B D^2+2 D E^2=B E^2+D F^2 .
$$

\section*{Доказательство через теорему косинусов}

Это просто применение теоремы косинусов. Мы знаем, что углы \( ADC \) и \( DCF \) являются дополнительными, поэтому их косинусы противоположны по знаку. Тогда теорема косинусов даёт нам:
\[ AC^2 = AD^2 + CD^2 - 2 \cdot AD \cdot CD \cdot \cos(ADC). \]
С другой стороны,
\[ DF^2 = AD^2 + CD^2 + 2 \cdot AD \cdot DC \cdot \cos(ADC). \]
Сложив эти равенства, мы получим свойство параллелограммов, использованное Эйлером.

Выразим из обоих уравнений \( DF^2 \), приравняем их и добавим \( AC^2 \):
\[ 2AD^2 + 2CD^2 = 2BD^2 + 2DE^2 + AC^2 - BE^2. \]

У нас остался ещё один параллелограмм \( ABCE \), для которого известно:
\[ 2AB^2 + 2BC^2 = AC^2 + BE^2. \]
Прибавим это к предыдущему уравнению:
\[ 2AD^2 + 2CD^2 + 2AB^2 + 2BC^2 = 2BD^2 + 2DE^2 + 2AC^2. \]

Разделим на два и немного преобразуем:
\[ AB^2 + BC^2 + CD^2 + DA^2 = AC^2 + BD^2 + DE^2, \]
что и требовалось доказать. \quad \textit{Q.E.D.}

\section*{Следствия Эйлера}

Эйлер приводит четыре следствия. Первые три довольно стандартны:

\begin{enumerate}
    \item \textbf{Если четырёхугольник — параллелограмм}, то отрезок \( DE \) вырождается, и сумма квадратов сторон в точности равна сумме квадратов диагоналей. Впрочем, это не совсем честно, поскольку именно это свойство параллелограммов Эйлер использовал в доказательстве самой теоремы.

    \item \textbf{Сумма квадратов сторон любого четырёхугольника всегда больше или равна сумме квадратов его диагоналей}, причём равенство достигается только для параллелограмма.

    \item \textbf{В третьем следствии} Эйлер делит диагональ \( AC \) пополам в точке \( P \), а диагональ \( BD \) — в точке \( Q \), и проводит отрезок \( QP \). Затем он показывает, что длина этого отрезка равна половине \( DE \), а его квадрат составляет четверть \( DE^2 \). Это удобно, так как позволяет избежать использования вспомогательной точки \( E \), заменив её двумя более естественными точками \( P \) и \( Q \), связанными с исходным четырёхугольником.
\end{enumerate}

Подставив \( 4PQ^2 \) вместо \( DE^2 \) в исходную теорему, мы приходим к \textbf{четвёртому и самому интересному следствию Эйлера}, проиллюстрированному на Рисунке 8.